%%%%%%%%%%%%%%%%%%%%%%%%%%%%%%%%%%%%%%%%%%%%%%%%%%%%%%%%%%%%%%%%%%%%%%%%%
%% Linear Algebra -- course project report template, Autumn 2026
%% Ukrainian Catholic University, Faculty of Applied Sciences,
%% Data Science Master Programme
%%
%% The title page is adapted from a template by WikiBooks, distributed
%% by LaTeXTemplates.com under CC BY-NC-SA 3.0.
%%
%% HOW TO USE THIS TEMPLATE
%%   1. Fill in the block marked "FILL IN" below: the stage of the report,
%%      the title, the authors, the project card and the links.
%%   2. One file serves all three milestones. Keep working in the same
%%      file: interim report 1 grows into interim report 2, and that grows
%%      into the final report.
%%   3. Replace the grey guidance boxes with your own text. Before you
%%      submit, set \showguidefalse so that any box you forgot disappears.
%%   4. Do NOT change the font size, the paper size or the margins.
%%      Page limits are counted in this format.
%%   5. Compile with pdfLaTeX (the Overleaf default). The folder must
%%      contain UCU-Apps.png (the logo); the example figure uses an image
%%      shipped with every TeX distribution.
%%
%% PAGE LIMITS (counted automatically; see the end of the log file)
%%   Interim reports 1 and 2: 2 pages, project card included.
%%   Final report: at most 10 pages from the abstract to the end of the
%%   Conclusion, figures and tables included. The title page, the
%%   references and the appendices are not counted.
%%%%%%%%%%%%%%%%%%%%%%%%%%%%%%%%%%%%%%%%%%%%%%%%%%%%%%%%%%%%%%%%%%%%%%%%%

\documentclass[11pt,a4paper]{article}

%----------------------------------------------------------------------
%  Page layout: do not change
%----------------------------------------------------------------------
\usepackage[a4paper,left=20mm,right=20mm,top=20mm,bottom=25mm,
            headheight=14pt,marginparwidth=16mm]{geometry}

%----------------------------------------------------------------------
%  Packages
%----------------------------------------------------------------------
\usepackage[english]{babel}
% T1 encoding only together with the scalable Latin Modern fonts
% (standard on Overleaf); otherwise fall back to the default encoding.
\IfFileExists{lmodern.sty}{\usepackage[T1]{fontenc}\usepackage{lmodern}}{}
\usepackage{microtype}
\usepackage{amsmath,amssymb,amsthm}
\usepackage{graphicx}
\usepackage{booktabs,tabularx}
\usepackage{caption,subcaption}
\usepackage[dvipsnames]{xcolor}
\usepackage{fancyhdr}
\usepackage{refcount}
\usepackage{listings}
\usepackage[ruled,vlined,linesnumbered]{algorithm2e}
% Margin notes are useful while drafting. Add the option "disable"
% before you submit, or simply delete your \todo{} commands.
\usepackage[colorinlistoftodos,textsize=scriptsize]{todonotes}
\usepackage[colorlinks=true,linkcolor=MidnightBlue,citecolor=MidnightBlue,
            urlcolor=MidnightBlue]{hyperref}   % load hyperref last

%%%%%%%%%%%%%%%%%%%%%%%%%%%%%%%%%%%%%%%%%%%%%%%%%%%%%%%%%%%%%%%%%%%%%%%%%
%%  FILL IN
%%%%%%%%%%%%%%%%%%%%%%%%%%%%%%%%%%%%%%%%%%%%%%%%%%%%%%%%%%%%%%%%%%%%%%%%%

% Stage of the report: 1 = interim report 1 (scoping memo),
%                      2 = interim report 2 (working baseline),
%                      3 = final report
\newcommand{\stage}{1}

\newcommand{\reporttitle}{Visual odometry for drone navigation}
\newcommand{\shorttitle}{Visual odometry for drone navigation}
\newcommand{\reportauthors}{Bohdan \textsc{Prokhorov}\\
  Markiian \textsc{Yatsyshyn}\\
  Vladyslav \textsc{Hryhola}}

% Project card: one line each. Keep it consistent with your registration;
% if something has changed, update it here and say so in the report.
\newcommand{\carddomain}{Computer vision and navigation}
\newcommand{\cardmethod}{SVD and low-rank approximation}
\newcommand{\cardobject}{The $N\times 9$ matrix $A$ from $N\ge 8$ point
  matches; the essential matrix $E$ reshaped from its null vector estimate,
  projected to singular values $(s,s,0)$}
\newcommand{\cardalgorithms}{Eight-point algorithm with and without Hartley
  normalisation; the same inside RANSAC (own implementations, library SVD)}
\newcommand{\carddata}{UZH-FPV Drone Racing Dataset~\cite{UZHFPV}, monocular image sequences}
\newcommand{\cardcode}{\url{https://github.com/your-team/your-repo}}
\newcommand{\cardvideo}{\url{https://link-to-your-recording}} % final only

% Guidance boxes: set \showguidefalse before you submit.
\newif\ifshowguide
\showguidefalse

%%%%%%%%%%%%%%%%%%%%%%%%%%%%%%%%%%%%%%%%%%%%%%%%%%%%%%%%%%%%%%%%%%%%%%%%%
%%  Template machinery: no need to read further than the next block
%%%%%%%%%%%%%%%%%%%%%%%%%%%%%%%%%%%%%%%%%%%%%%%%%%%%%%%%%%%%%%%%%%%%%%%%%

% Notation used in the course. Add your own macros here.
\newcommand{\bR}{\mathbb{R}}
\newcommand{\bC}{\mathbb{C}}
\newcommand{\ba}{\mathbf{a}}
\newcommand{\bb}{\mathbf{b}}
\newcommand{\bu}{\mathbf{u}}
\newcommand{\bv}{\mathbf{v}}
\newcommand{\bx}{\mathbf{x}}
\newcommand{\by}{\mathbf{y}}
\newcommand{\bz}{\mathbf{z}}
\DeclareMathOperator{\rank}{rank}
\DeclareMathOperator{\cond}{cond}
\DeclareMathOperator{\diag}{diag}
\DeclareMathOperator{\trace}{trace}
\DeclareMathOperator{\nnz}{nnz}
\DeclareMathOperator*{\argmin}{argmin}

% Theorem-like environments, numbered within sections
\theoremstyle{plain}
\newtheorem{theorem}{Theorem}[section]
\newtheorem{proposition}[theorem]{Proposition}
\newtheorem{lemma}[theorem]{Lemma}
\theoremstyle{definition}
\newtheorem{definition}[theorem]{Definition}
\newtheorem{example}[theorem]{Example}
\theoremstyle{remark}
\newtheorem{remark}[theorem]{Remark}

% Code listings
\lstset{language=Python, basicstyle=\ttfamily\small, columns=fullflexible,
  keepspaces=true, frame=lines, numbers=left, numberstyle=\tiny\color{gray},
  commentstyle=\color{gray}, keywordstyle=\color{MidnightBlue}\bfseries,
  showstringspaces=false, breaklines=true}

% Guidance boxes
\newcommand{\guide}[1]{\ifshowguide\par\smallskip\noindent
  \colorbox{gray!12}{\parbox{\dimexpr\linewidth-2\fboxsep}{\small\itshape
  \textbf{Guidance.} #1}}\par\smallskip\fi}

% Stage names
\newcommand{\stagename}{\ifcase\stage\or Interim report 1: scoping memo\or
  Interim report 2: working baseline\or Final report\fi}

% Running heads
\pagestyle{fancy}
\fancyhf{}
\fancyhead[L]{\small\shorttitle}
\fancyhead[R]{\small\stagename}
\fancyfoot[C]{\small\thepage}
\renewcommand{\headrulewidth}{0.4pt}

% Project card
\newcommand{\projectcard}{%
  \noindent
  \begin{tabularx}{\linewidth}{@{}>{\raggedright\bfseries}p{3.0cm}
                                  >{\raggedright\arraybackslash}X@{}}
    \toprule
    Domain        & \carddomain \\
    Method family & \cardmethod \\
    LA object     & \cardobject \\
    Algorithms    & \cardalgorithms \\
    Data          & \carddata \\
    \ifnum\stage>1 Code & \cardcode \\ \fi
    \ifnum\stage=3 Recording & \cardvideo \\ \fi
    \bottomrule
  \end{tabularx}\par}

% Page counting: the main body runs from \bodystart to \bodyend
\newcommand{\pagelimit}{\ifnum\stage=3 10\else 2\fi}
\newcommand{\bodystart}{\label{body-start}}
\newcommand{\bodyend}{\label{body-end}}
\AtEndDocument{%
  \refused{body-start}\refused{body-end}%
  \edef\bodypages{\number\numexpr\getpagerefnumber{body-end}
    -\getpagerefnumber{body-start}+1\relax}%
  \typeout{}%
  \typeout{*** Main body: \bodypages\space page(s); limit \pagelimit.}%
  \typeout{}%
  \ifnum\bodypages>\pagelimit\relax
    \PackageWarningNoLine{LA-template}{Main body has \bodypages\space pages,
      above the limit of \pagelimit}%
  \fi}

%%%%%%%%%%%%%%%%%%%%%%%%%%%%%%%%%%%%%%%%%%%%%%%%%%%%%%%%%%%%%%%%%%%%%%%%%
\begin{document}
%%%%%%%%%%%%%%%%%%%%%%%%%%%%%%%%%%%%%%%%%%%%%%%%%%%%%%%%%%%%%%%%%%%%%%%%%

\ifnum\stage=3
%----------------------------------------------------------------------
%  Final report: title page (not counted)
%----------------------------------------------------------------------
\begin{titlepage}
  \centering
  {\large\scshape Ukrainian Catholic University\par}
  {\scshape Faculty of Applied Sciences\quad$\cdot$\quad
    Data Science Master Programme\par}
  \vspace{1.6cm}
  \rule{\linewidth}{0.5mm}\par\vspace{0.4cm}
  {\LARGE\bfseries \reporttitle\par}
  \vspace{0.3cm}
  {\large Linear Algebra: course project, final report\par}
  \vspace{0.2cm}
  \rule{\linewidth}{0.5mm}\par
  \vspace{1.2cm}
  {\large\emph{Authors:}\par\smallskip \reportauthors\par}
  \vspace{1.2cm}
  \projectcard
  \vfill
    \IfFileExists{UCU-Apps.png}{\includegraphics[height=4cm]{UCU-Apps.png}\par}{}
  \vspace{0.2cm}
  {\large \today\par}
\end{titlepage}

%----------------------------------------------------------------------
%  Abstract: opens the first page of the main body
%----------------------------------------------------------------------
\bodystart
\begin{abstract}
\guide{Five to eight sentences: the problem, the linear-algebra object,
the methods you compared, the main quantitative result, and the main
limitation. Write it last.}
We compare two algorithms for completing a partially observed ratings
matrix by a matrix of low rank~\cite[Sec.~III.5]{Strang}.
\end{abstract}

\else
%----------------------------------------------------------------------
%  Interim reports: compact header and project card (counted)
%----------------------------------------------------------------------
\bodystart
\begin{center}
  {\Large\bfseries \reporttitle\par}
  \medskip
  {\large \stagename\par}
  \smallskip
  \begin{tabular}{c}\reportauthors\end{tabular}\par
  \smallskip
  {\small Linear Algebra, UCU Data Science Master Programme, \today\par}
\end{center}
\projectcard
\medskip
\fi


%%%%%%%%%%%%%%%%%%%%%%%%%%%%%%%%%%%%%%%%%%%%%%%%%%%%%%%%%%%%%%%%%%%%%%%%%
\section{Problem and aim}
\label{sec:problem}

\guide{Criterion A. What goes in, what comes out, and why it matters in
practice. Who uses a solution to this problem, and what goes wrong
without a good one? Give a short overview of the approaches in the
literature and say which ones you chose and why.
\emph{Interim report 1}: half a page.}

The input is a video from an FPV racing drone with known camera calibration, with features matched between consecutive frames, which gives $N$ point correspondences per frame pair. The output is the relative rotation $R$ and the direction of the translation $\mathbf{t}$ for each pair, chained into the camera trajectory. Such camera-based motion estimation lets drones, robots and cars navigate where GPS is unavailable or unreliable~\cite{ScaramuzzaI}. Racing drones are a hard case, because fast motion blurs the images and produces many wrong matches~\cite{UZHFPV}, and small errors accumulate into drift. With one camera the scale of $\mathbf{t}$ is unknown.

For every frame pair we measure the rotation error $e_R=\arccos\frac{\trace(R^{\top}R^{*})-1}{2}$ and the translation-direction error $e_t=\angle(\mathbf{t},\mathbf{t}^{*})$ against the truth $(R^{*},\mathbf{t}^{*})$, and the running time. The project is successful if (i) on synthetic data with $1$ pixel noise the normalized eight-point algorithm gives $e_R<1^{\circ}$ and $e_t<5^{\circ}$, and (ii) on the UZH-FPV sequences (one Easy and one Hard, with public ground truth) Hartley normalization lowers the median errors compared with the plain algorithm, and our RANSAC lowers the $90$th percentile errors further.

%%%%
\section{Linear-algebra formulation}
\label{sec:formulation}

\guide{Criterion A. State the problem precisely as a linear-algebra
problem: the matrix (its size, its entries, its structure: sparse,
symmetric, positive definite, low rank), the unknowns, and what you need
to do with the matrix. Explain why this structure suggests your methods.
\emph{Interim report 1}: this section is the core of the memo.}

A 3D point seen in two frames gives calibrated image points $\bx_1,\bx_2\in\bR^3$ (homogeneous coordinates), which satisfy the epipolar constraint
\begin{equation}\label{eq:epi}
  \bx_2^{\top}E\,\bx_1=0,
  \qquad
  E=[\mathbf{t}]_{\times}R,
\end{equation}
where $[\mathbf{t}]_{\times}$ is the matrix of the cross product with $\mathbf{t}$. The essential matrix $E$ has singular values $(s,s,0)$ with $s=\|\mathbf{t}\|$~\cite[Ch.~9]{HartleyZisserman}, and it is defined only up to scale, so we take $\|\mathbf{t}\|=1$. Equation~\eqref{eq:epi} is linear in the nine entries $\mathbf{e}\in\bR^9$ of $E$, so $N$ matches give
\begin{equation}\label{eq:Ae}
  A\mathbf{e}=0,
  \qquad
  A\in\bR^{N\times 9},
  \quad
  \mathbf{a}_i^{\top}=(x_2x_1,\,x_2y_1,\,x_2,\,y_2x_1,\,y_2y_1,\,y_2,\,x_1,\,y_1,\,1)_i .
\end{equation}
The central object is this dense matrix with $N\ge 8$ (in practice hundreds of rows). For nondegenerate, noiseless correspondences $\rank A=8$ and $\mathbf{e}$ spans its null space. With noise we solve
\begin{equation}\label{eq:hls}
  \min_{\|\mathbf{e}\|_2=1}\|A\mathbf{e}\|_2,
\end{equation}
SVD directly solves this constrained minimization: choose the last right singular vector $\mathbf{v}_9$~\cite[App.~A5]{HartleyZisserman}. For $N=8$, use the full $9\times9$ right-singular basis; a reduced SVD omits the null vector. Homogeneity does not rule out QR-based methods. Coordinate scaling can harm numerical accuracy, especially in pixels. However, for $N\ge9$ exact rank-eight data already gives $\cond_2(A)=\infty$. We must distinguish scaling of the nonzero spectrum from separation of the smallest singular direction.

%%%%
\section{Methods}
\label{sec:methods}

\guide{Criteria B and C. At most two algorithms. For each: the key idea,
explained in your own words; the central claims, proved or sketched
(say which, and cite the source of any proof you do not give); the cost
in flops and memory, as a function of the problem size; and the
conditions under which it fails or becomes slow.
\emph{Interim report 1}: the candidate approaches, their pros and cons,
and their predicted cost.}

After solving~\eqref{eq:hls} and undoing any Hartley normalization, let $\hat E$ be the matrix estimate in calibrated camera coordinates. Project it onto the essential matrices: if $\hat E=U\diag(\sigma_1,\sigma_2,\sigma_3)V^{\top}$, then $E=U\diag(s,s,0)V^{\top}$ with $s=\frac12(\sigma_1+\sigma_2)$ is the closest one in the Frobenius norm. Then, with $W=\bigl(\begin{smallmatrix}0&-1&0\\1&0&0\\0&0&1\end{smallmatrix}\bigr)$, $\mathbf{u}_3$ the last column of $U$ and $\det U=\det V=1$, we factor
\begin{equation}\label{eq:decomp}
  (R,\mathbf{t})\in\bigl\{UWV^{\top},\;UW^{\top}V^{\top}\bigr\}\times\{\mathbf{u}_3,-\mathbf{u}_3\}.
\end{equation}
Of these four candidates we keep the one for which the triangulated points (again an SVD problem) are in front of both cameras.

Table~\ref{tab:approaches} compares three variants. The SVD of an $N\times n$ matrix with $N\ge n$ costs $O(Nn^2)$ flops, hence $O(N)$ for $n=9$. The plain eight-point algorithm solves~\eqref{eq:hls} directly~\cite{Hartley97}. Hartley normalization translates each frame's points to centroid zero and scales their mean distance to $\sqrt2$, using $T_1,T_2$. Solve for a matrix $M$ in these coordinates, undo the transforms as $\hat E=T_2^{\top}MT_1$, then enforce the essential constraint above. The transformed matrix $M$ need not have equal nonzero singular values. For the pixel-coordinate experiment, instead estimate a rank-two fundamental matrix $F$ (without requiring equal singular values), denormalize it, and form $\hat E=K_2^{\top}FK_1$ using the camera calibration matrices before essential projection.

RANSAC takes 8 random matches, estimates $E$, counts inliers using a residual threshold, and refits on the largest inlier set~\cite{HartleyZisserman,ScaramuzzaII}. The residual and threshold remain to be agreed; raw $|\bx_2^{\top}E\bx_1|$ depends on the scale of $E$. For inlier ratio $w$ and confidence $p$, $K=\lceil\ln(1-p)/\ln(1-w^8)\rceil$, about 1200 for $w=0.5$, $p=0.99$. Each iteration costs $O(N)$ for scoring plus constant $c_8$ for the $8\times9$ solve.

For $N\ge9$, a small gap between $\sigma_8$ and $\sigma_9$ indicates sensitivity of the last singular direction, not nonuniqueness unless they are equal (apart from sign). Exact planar degeneracy can give $\sigma_8=\sigma_9=0$, so their ratio is undefined. We will test planar scenes and small baselines to assess these diagnostics; no universal detection threshold is assumed.

\begin{table}[ht]
  \centering
  \caption{Approaches, predicted cost for $N$ matches, and what is likely to go wrong.}
  \label{tab:approaches}
  \begin{tabularx}{\linewidth}{@{}lllX@{}}
    \toprule
    Approach & Flops & Memory & Likely problems \\
    \midrule
    Plain 8-point      & $O(Nn^2)$, $n=9$ & $O(N)$ & poor coordinate scaling; noise; wrong matches \\
    Normalized 8-point & $O(Nn^2)+O(N)$ & $O(N)$ & wrong matches; planar degeneracy; small baseline \\
    RANSAC + normalized & $K\,(c_8+O(N))$ & $O(N)$ & $K$ grows like $w^{-8}$ for few inliers; random result \\
    \bottomrule
  \end{tabularx}
\end{table}

\ifnum\stage>1
\section{Implementation and experiments}
\label{sec:experiments}

\guide{Criterion D. What you implemented yourselves, and what you called
from a library. The data, the benchmark and the metrics. Then the
results: accuracy, the condition numbers of the matrices you actually
factor, runtime against problem size next to the predicted cost, and the
size or condition number at which a method stops working.
\emph{Interim report 2}: a first implementation running on small data
and one comparison with a library routine, with one figure.}

\subsection{Setup}

Describe the hardware and the software versions briefly, so that a
reader can judge the timings. Every figure and table must be
reproducible from the \texttt{README} of your repository.

\subsection{Results}

Figure~\ref{fig:example} shows how to include a figure, and
Table~\ref{tab:example} how to report a comparison. Save plots from
Python as vector PDF, with \texttt{plt.savefig("fig.pdf")} rather than
as PNG. Label both axes, with units, and use a logarithmic axis when
values span several orders of magnitude.

\begin{figure}[ht]
  \centering
  \begin{subfigure}{0.45\linewidth}
    \includegraphics[width=\linewidth]{example-image-a}
    \caption{Runtime against $n$.}
  \end{subfigure}\hfill
  \begin{subfigure}{0.45\linewidth}
    \includegraphics[width=\linewidth]{example-image-b}
    \caption{Relative error against $\cond(A)$.}
  \end{subfigure}
  \caption{A figure with two panels. The caption says what the reader
    should see, not only what is plotted.}
  \label{fig:example}
\end{figure}

\begin{table}[ht]
  \centering
  \caption{Comparison of methods. Report the mean and the spread over
    several runs, and say how many runs.}
  \label{tab:example}
  \begin{tabular}{lrrr}
    \toprule
    Method & $n$ & Time (s) & Relative error \\
    \midrule
    Algorithm 1 (ours)      & \dots & \dots & \dots \\
    Algorithm 2 (ours)      & \dots & \dots & \dots \\
    Library routine         & \dots & \dots & \dots \\
    \bottomrule
  \end{tabular}
\end{table}

%%%%%%%%%%%%%%%%%%%%%%%%%%%%%%%%%%%%%%%%%%%%%%%%%%%%%%%%%%%%%%%%%%%%%
\fi
\ifnum\stage=3
\section{Discussion and limitations}
\label{sec:discussion}

\guide{Criteria C and D. Did the experiments confirm the predicted cost
and failure modes? Explain every surprise. What did not work, and why?
An honest explanation of a failure earns credit; a clean result with no
explanation does not.}

%%%%%%%%%%%%%%%%%%%%%%%%%%%%%%%%%%%%%%%%%%%%%%%%%%%%%%%%%%%%%%%%%%%%%%%%%
\section{Conclusion}
\label{sec:conclusion}

\guide{What you achieved, the main difficulties, and what you would do
next. Half a page is enough.}

\fi

\ifnum\stage<3
%%%%%%%%%%%%%%%%%%%%%%%%%%%%%%%%%%%%%%%%%%%%%%%%%%%%%%%%%%%%%%%%%%%%%%%%%
\section{Chosen methods and plan}
\label{sec:plan}

\guide{Interim reports only; this section disappears from the final
report. \emph{Interim report 1}: the planned steps and who does what.
\emph{Interim report 2}: what is not working yet, any change of plan,
and the outline of the final report.}

We compare plain and normalized eight-point estimation, adding RANSAC to test outlier rejection. We implement $A$, normalization, essential projection and decomposition, triangulation, RANSAC and the error metrics, using a library SVD. OpenCV's five-point estimator is the independent baseline, evaluated on the same matches with the same pose-error metrics. A common feature interface will support classical and neural-network detectors and matchers.

Until 30 Nov (interim report 2), we plan camera-model checks and fisheye undistortion if needed, the feature pipeline, synthetic data with known poses and noise, the metrics and plain eight-point estimation. Until 21 Dec (final report), we plan normalization, RANSAC, the library comparison, synthetic noise sweeps including planar scenes and small baselines, UZH-FPV Easy/Hard runs with different matchers, runtime scaling, and the report. Work packages are data/pipeline, features, synthetic data/metrics, eight-point estimation, RANSAC, experiments and writing; owners are not yet decided.

AI assistance: Codex helped review and correct mathematics, references and wording.

\fi

\bodyend

%%%%%%%%%%%%%%%%%%%%%%%%%%%%%%%%%%%%%%%%%%%%%%%%%%%%%%%%%%%%%%%%%%%%%%%%%
%%  References (not counted)
%%%%%%%%%%%%%%%%%%%%%%%%%%%%%%%%%%%%%%%%%%%%%%%%%%%%%%%%%%%%%%%%%%%%%%%%%
% Cite only sources you have actually read. Give a DOI or a link where
% one exists. Cite the datasets and the main software you used.
\begin{thebibliography}{9}

\bibitem{Hartley97}
R.~I.~Hartley,
In defense of the eight-point algorithm,
\emph{IEEE Transactions on Pattern Analysis and Machine Intelligence}
\textbf{19}(6), 580--593, 1997.
\url{https://www.cs.ucf.edu/courses/cap6938-02/refs/hartley.pdf}.

\bibitem{HartleyZisserman}
R.~Hartley and A.~Zisserman,
\emph{Multiple View Geometry in Computer Vision}, 2nd ed.,
Cambridge University Press, 2004.
\url{https://www.robots.ox.ac.uk/~vgg/hzbook/}.

\bibitem{ScaramuzzaI}
D.~Scaramuzza and F.~Fraundorfer,
Visual odometry [Tutorial],
\emph{IEEE Robotics \& Automation Magazine}
\textbf{18}(4), 80--92, 2011.
\url{https://doi.org/10.1109/MRA.2011.943233}.

\bibitem{ScaramuzzaII}
F.~Fraundorfer and D.~Scaramuzza,
Visual odometry: Part II: Matching, robustness, optimization, and applications,
\emph{IEEE Robotics \& Automation Magazine}
\textbf{19}(2), 78--90, 2012.
\url{https://doi.org/10.1109/MRA.2012.2182810}.

\bibitem{UZHFPV}
J.~Delmerico, T.~Cieslewski, H.~Rebecq, M.~Faessler and D.~Scaramuzza,
Are we ready for autonomous drone racing? The UZH-FPV drone racing dataset,
\emph{IEEE International Conference on Robotics and Automation (ICRA)}, 2019.
\url{https://fpv.ifi.uzh.ch/}.

\end{thebibliography}

%%%%%%%%%%%%%%%%%%%%%%%%%%%%%%%%%%%%%%%%%%%%%%%%%%%%%%%%%%%%%%%%%%%%%%%%%
%%  Appendices (not counted)
%%%%%%%%%%%%%%%%%%%%%%%%%%%%%%%%%%%%%%%%%%%%%%%%%%%%%%%%%%%%%%%%%%%%%%%%%
\ifnum\stage=3
\appendix

\section{Contribution statement}

\guide{Required. Who did what, agreed by all members. Staff may adjust
an individual grade if contributions are clearly unequal.}

\noindent
\begin{tabularx}{\linewidth}{@{}lX@{}}
  \toprule
  Member & Contribution \\
  \midrule
  A.~A.~Lovelace & \dots \\
  C.~Babbage     & \dots \\
  \bottomrule
\end{tabularx}

\section{AI-use statement}

\guide{Required. Which AI tools you used and for what: code, language
editing, finding literature. Every member must be able to explain every
equation and every line of the core code. If you used none, say so.}

\section{Additional results}

\guide{Optional. Proofs, extra figures and tables that support the main
text. Reviewers are not required to read the appendices, so everything
needed to follow the report belongs in the main body.}
\fi

 \end{document}
